\begin{definition}[Jordan Block]
    A Jordan block $J_{\lambda, k}$ of eigenvealue $\lambda$ and index/size $k$ is a $k\times k$ upper triangular matrix such that$\cdots$ 
\end{definition}

\begin{theorem}
Let $J=J_{\lambda, k}\in M_k(\mathbb{C}),$ then:
\begin{enumerate}
    \item $p_J(x)=(\lambda-x)^k.$
    \item $a_J(\lambda)=k,$ and $g_J(\lambda)=1.$
    \item For $1\le i\le k,$ $\rank((J-\lambda I_k)^i)=k-i$ and $\nullity((J-\lambda I_k)^i)= i.$
    \item $\nullity((J-\lambda I_k)^i)-\nullity((J-\lambda I_k)^{i-1}) =
\begin{cases}
1, & \text{if } 1 \leq i \leq k \\
0, & \text{if } i \geq k+1
\end{cases}.$
    \item $m_J(x)=(x-\lambda)^k.$
\end{enumerate}
\label{JB_property}
\end{theorem}



\begin{pfbox}
1,2 is trivial by inspection. For 3 and 4, let $N_k=J-\lambda I_k$. The key observation is that the diagonal of $1'$s is shifted by one unit to the right after multiplying $N_k$ itself.
For 5, since $N_k^i\neq \mathbf{O}$ for $i < k$, the minimal polynomial of $J$ must be the characteristic polynomial of $J$.
\end{pfbox}


\begin{remark}
2 and 5 shows that $J$ is diagonalizable only when $J$ is of size 1.
\end{remark}


\begin{definition}{Direct sum of matrices}
  Let $A\in M_m(F),\,B\in M_n(F)$, then the direct sum $A\oplus B$ is a block diagonal matrix in $M_{m+n}(F)$ such that
  \[A\oplus B = \begin{pmatrix}
    A & \mathbf{O}\\
    \mathbf{O} & B.
  \end{pmatrix}\]
\end{definition}

\begin{exercise}
    Write down the definition of a direct sum of two subspaces. They are two different concepts.
\end{exercise}

\begin{theorem}
     Let $A\in M_m(F),\,B\in M_n(F)$, then:
     \begin{enumerate}
        \item $\rank(A\oplus B) = \rank(A)+\rank(B)$
        \item $\nullity(A\oplus B) = \nullity(A)+\nullity(B)$
        \item $p_{A\oplus B}(x)=p_A(x)\cdot p_B(x)$
        \item $a_{A\oplus B}(\lambda)=a_A(\lambda)+a_B(\lambda)$
        \item $g_{A\oplus B}(\lambda)=g_A(\lambda)+g_B(\lambda)$
        \item $m_{A\oplus B}(x)=\text{lcm}\{m_A(x), m_B(x)\}$
     \end{enumerate}
\end{theorem}


\begin{pfbox}
    1--5 are left as an exercise. 
    For 6, observe that for any polynomial $f(x)\in F[x]$, one has $f(A\oplus B)=f(A)\oplus f(B)$.
    Let $m_{A\oplus B}=f(x),\, f(A\oplus B)=\mathbf{0}$ by the definition of minimal polynomial, which forces $f(x)$ is a common multiple of $m_A(x)$ and $m_B(x).$ 
    For the "least" part, we provide two methods. Suppose $g(x)$ is another common multiple, we want to show $f$ divides $g$.
    Since $m_A(x)\mid g(x)$ and $m_B(x)\mid g(x),\, g(A\oplus B) = \mathbf{0}$ impies $f\mid g.$ \\

    Alternatively, conduct the division $g(x)=f(x)\cdot q(x)+r(x).$ It suffices to show $r(x)\equiv 0.$
    Suppose, by contradiction, $deg(r)<deg(f).$ However,$r(A)=r(B)=\mathbf{0}$ violates the minimality.
\end{pfbox}




\begin{definition}[JCF]
    A matrix $J$ is in its Jordan Canonical Form if it is a direct sum of Jordan blocks:
    \[J=J_{\lambda_1,k_1}\oplus J_{\lambda_2,k_2}\oplus\cdots\oplus J_{\lambda_n,k_n}.\]
\end{definition}

\begin{corollary}
    \begin{enumerate}
        \item Any diagonal matrix is exactly in its JCF$\cdots$
        \item Write down the $p_J(x)$ and $m_J(x)$ in terms of $p_{J_{\lambda_i, k_i}}(x)$ and $m_{J_{\lambda_i, k_i}}(x)$.
    \end{enumerate}
\end{corollary}


\begin{theorem}
    Let $J$ in its JCF and fix $\lambda\in Spec(J),$ then:
    \begin{enumerate}
        \item \# of Jordan blocks of eigenvalue $\lambda = g_J(\lambda).$
        \item \# of Jordan blocks of eigenvalue $\lambda$ of index(size) at least $i = \nullity((J-\lambda I_k)^i) - \nullity((J-\lambda I_k)^{i-1}).$
    \end{enumerate}
\end{theorem}


\begin{pfbox}
    Let $J=J_{\lambda_1,k_1}\oplus J_{\lambda_2,k_2}\oplus\cdots\oplus J_{\lambda_n,k_n}.$ Suppose $\lambda_1=\lambda_2=\cdots\lambda_l=\lambda,\, l\le n,\, \lambda\notin\{\lambda_{l+1},\lambda_{l+2},\cdots\lambda_{n}\}.$
    Let $N_k=J-\lambda I_k,$ recall (3) and (4) in theorem~\eqref{JB_property},
    \begin{enumerate}
        \item $g_J(\lambda)=\sum_{i=1}^{n}\nullity(N_{\lambda_i, k_i}) = \underbrace{1+1+\cdots+1}_{l \text{ times}} + \underbrace{0+0+\cdots+0}_{\lambda\ne\lambda_i,\, i\ge l}=l$,
        which equals \# of Jordan blocks of eigenvalue $\lambda.$
        \item Because $N_k=N_{k_1}\oplus N_{k_2}\oplus\cdots\oplus N_{k_l}\oplus T,$ where $T$ is some direct sum of Jordan blocks of eigenvalue that is not $\lambda.$
        Fix $i,\,\nullity(N^i_k)-\nullity(N^{i-1}_k)=\sum_{j=1}^{l}\left(\nullity(N^i_{k_j})-\nullity(N^{i-1}_{k_j})\right)=\textcolor{darkred}{\#\{1\le j\le l:i\le k_j\}}$, which equals
        \# of Jordan blocks of eigenvalue $\lambda$ of index(size) at least $i$.
    \end{enumerate}
\label{determine_JCF_block_size}
\end{pfbox}


Combine the definition of JCF and the theroem~\eqref{determine_JCF_block_size}, we can define the JCF of a general matrix $A\in M_n(\mathbb{C}).$

\begin{definition}[JCF of a general matrix]
    Let $A\in M_n(\mathbb{C}),$ the JCF of $A$ is a matrix of $J_A$ s.t.
    \begin{enumerate}
        \item $J_A$ is a direct sum of Jordan blocks.
        \item \# of Jordan blocks of eigenvalue $\lambda = g_A(\lambda).$
        \item \# of Jordan blocks of eigenvalue $\lambda$ of index(size) at least $i = \nullity\left((A -\lambda I_k)^i\right) - \nullity(\left(A -\lambda I_k)^{i-1}\right).$
    \end{enumerate}
\end{definition}

\begin{gbox}
\begin{remark}
How to find the JCF of $A ?$
\begin{enumerate}
    \item Find spectrum of $A$ via $p_A(x).$
    \item Find $g_A(\lambda)$ to determine \# of Jordan blocks.
    \item Find $r_i=\nullity\left((A -\lambda I_k)^i\right) - \nullity(\left(A -\lambda I_k)^{i-1}\right)$ to determine individual size of a Jordan block.
\end{enumerate}
\end{remark}
\end{gbox}




Let $T:V\to V$ be a linear operator over a finite-dimensional complex vector space $V$.
In the context of diagonalization, we aim to find an eigenbasis $\eta$ such that $[T]^\eta_\eta$ is a diagonal matrix.
Likewise, our goal now is to find a basis $\gamma$ such that $[T]^\gamma_\gamma$ is in its JCF.
We first look at a part of matrix corresponding to an eigenvalue $\lambda$. Suppose there exists such a part of basis $\beta=\{\mathbf{v}_1, \mathbf{v}_2,\cdots\mathbf{v}_k\}$, and we put each elements by columns to form a matrix $P$, then we have:
\begin{align*}
[T]^B_B P &= P [T]^\beta_\beta \\
T([\mathbf{v}_1, \mathbf{v}_2, \dots, \mathbf{v}_k]) &= [\mathbf{v}_1, \mathbf{v}_2, \dots, \mathbf{v}_k] 
\begin{pmatrix} 
\lambda & 1 & 0 & \cdots & 0 \\ 
0 & \lambda & 1 & \cdots & 0 \\ 
\vdots & \vdots & \ddots & \ddots & \vdots \\ 
0 & 0 & \cdots & \lambda & 1 \\ 
0 & 0 & \cdots & 0 & \lambda 
\end{pmatrix},
\end{align*}

where $B$ is the standard basis. Comparing the columns on both sides yields the following system of equations:

$$
\begin{aligned}
T(\mathbf{v}_1) &= \lambda \mathbf{v}_1 \\
T(\mathbf{v}_2) &= \mathbf{v}_1 + \lambda \mathbf{v}_2 \\
T(\mathbf{v}_3) &= \mathbf{v}_2 + \lambda \mathbf{v}_3 \\
&\ \ \vdots \\
T(\mathbf{v}_k) &= \mathbf{v}_{k-1} + \lambda \mathbf{v}_k
\end{aligned}
\quad
\Longleftrightarrow
\quad
\begin{aligned}
(T - \lambda \operatorname{id})(\mathbf{v}_1) &= \mathbf{0} \\
(T - \lambda \operatorname{id})(\mathbf{v}_2) &= \mathbf{v}_1 \\
(T - \lambda \operatorname{id})(\mathbf{v}_3) &= \mathbf{v}_2 \\
&\ \ \vdots \\
(T - \lambda \operatorname{id})(\mathbf{v}_k) &= \mathbf{v}_{k-1}
\end{aligned}
$$

Observe that:
\begin{enumerate}
    \item $\mathbf{v}_1$ is the only eigenvector among $\beta.$
    \item Given $\mathbf{v}_k,$ we can generate a chain of $\beta$ by setting $\Big\{\mathbf{v}_k, (T-\lambda \operatorname{id})(\mathbf{v}_k), (T-\lambda \operatorname{id})(\mathbf{v}_{k-1}),\cdots(T-\lambda \operatorname{id})(\mathbf{v}_2)\Big\},$
    or equivalently,\\
     $\Big\{\mathbf{v}_k, (T-\lambda \operatorname{id})(\mathbf{v}_k), (T-\lambda \operatorname{id})^2(\mathbf{v}_{k}),\cdots,(T-\lambda \operatorname{id})^{k-1}(\mathbf{v}_{k})\Big\}.$
    \item $(T-\lambda \operatorname{id})^{k}(\mathbf{v}_{k})=(T-\lambda \operatorname{id})(\mathbf{v}_{1})=\mathbf{0},$ so $\mathbf{v}_k\in \ker(T-\lambda\operatorname{id})^k\setminus\ker(T-\lambda\operatorname{id})^{k-1} (\because \mathbf{v}_1\ne\mathbf{0}).$
\end{enumerate}

These motivates the definition of Jordan chain:
\begin{definition}[Jordan Chain]
    Let $T:V\to V$ be a linear operator and $\lambda\in Spec(T).$
    We say a set of non-zero vectors $\beta=\{\mathbf{v}_1, \mathbf{v}_2,\cdots\mathbf{v}_k\}$ is a Jordan chain of $T$ corresponding to the eigenvalue $\lambda$ of length $k$ if $\dots$,
    where $\mathbf{v}_k$ is the end vector and $\mathbf{v}_1$ is the initial vector.
\end{definition}


\begin{definition}[Generalized eigenspace]
    Let $T:V\to V$ be a linear operator and $\lambda$ be an eigenvalue of $T$.
    Define 
    \[K_\lambda\equiv\bigcup_{i\ge 1}\ker(T-\lambda)^i\]
    to be the generalized eigenspace of $T$ w.r.t eigenvalue $\lambda$. Any non-zero vector in $K_\lambda$ is called a generalized eigenvector.
\end{definition}

\begin{theorem}
    Suppose $m$ is the smallest positive integer such that $\ker(T-\lambda I)^m=\ker(T-\lambda I)^{m+1}.$
    Then 
    \begin{enumerate}
    \item we have a strict inclusions of subspaces:
    \[\{\mathbf{0}\}\subset\ker(T-\lambda I)\subset\ker(T-\lambda I)^2\subset\cdots\subset\ker(T-\lambda I)^m.\]
    \item $K_\lambda=\ker(T-\lambda I)^m.$
    \item Let $p\le m$ and $\mathbf{v}\in\ker(T-\lambda I)^p\setminus\ker(T-\lambda I)^{p-1}.$\\
    The set $\{\mathbf{v},(T-\lambda I)(\mathbf{v}),\cdots,(T-\lambda I)^{p-1}(\mathbf{v})\}$ forms a Jordan chain of $T$ w.r.t eigenvalue $\lambda$ of length $p.$
    \end{enumerate}
    \label{thm 7.4}
\end{theorem}


\begin{theorem}
    Let $\lambda$ be an eigenvalue of $T.$ If $\beta_1,\beta_2,\cdots\beta_q$ are Jordan chains of $\lambda$ such that their initial vectors are linearly independent, then 
    $\gamma_\lambda = \beta_1\cup\beta_2\cup\cdots\cup\beta_q$ is linearly independent.
\end{theorem}

\begin{pfbox}
Let $|\gamma_\lambda|=n,$ we aim to do induction on the size of union of Jordan chains, which are linearly independent.
If $n=1,$ there are nothing to prove. Suppose $n<k$ is true, then for $n=k,$ let $W = span(\gamma_\lambda),$ we want to show $dim(W)=k.$
Consider a restrction map $T-\lambda I|_W:W\to W$ and a set $\gamma'_\lambda=\beta'_1\cup\beta'_2\cup\cdots\beta'_q$ obtained by removing the end vector of each Jordan chain.
By induction hypothesis, $\gamma'_\lambda$ is linearly independent. Note that $\gamma_\lambda$ spans $W,\,(T-\lambda I)(\gamma_\lambda)=\gamma'_\lambda$ spans $(T-\lambda)(W),$ i.e. $im(T-\lambda|_W).$
So $\gamma'_\lambda$ forms a basis of $im(T-\lambda I_W)$ and $\rank(T-\lambda I)=k-q.$ Hence 
$$\underbrace{k}_{|\gamma_\lambda|}\ge dim(W)=\rank(T-\lambda|_W)+\nullity(T-\lambda I)\ge k-q + q = k,$$
where $\ker(T-\lambda I|_W)=\ker(T-\lambda I)\cap W$ has at least $q$ linearly independent eigenvector (initial vectors).
This forces $dim(W)=k.$
\end{pfbox}


\begin{theorem}
    There exists a basis $\gamma_\lambda$ of generalized eigenspace $K_\lambda$ such that $\gamma_\lambda$ is formed by a union of Jordan chains of the eigenvalue $\lambda$.
    \label{thm 7.6}
\end{theorem}

\begin{pfbox}
    We prove by the induction on the dimension of $K_\lambda.$ Let $n=dim(K_\lambda),\,n=1,$ there are nothing to prove.
    Now suppose $n<k$ such that $K_\lambda$ admits a basis formed by a union of Jordan chains, then for $n=k,$ consider a restriction $T-\lambda I|_{K_\lambda}: K_\lambda\to K_\lambda.$
    The proof strategy goes as following:
    \begin{enumerate}
        \item $ker(T-\lambda I)\supset\{\mathbf{0}\},$ by rank-nullity, $im(T-\lambda I|_{K_\lambda})\subset K_\lambda.$
        \item Claim: Let $u=im(T-\lambda I|_{K_\lambda}),\, u$ is the generalized eigenspace of $T|_u$ w.r.t eigenvalue $\lambda$. From this claim and induction hypothesis, there exist Jordan chains $\beta_1,\beta_2,\cdots,\beta_q$ such that their union, $\gamma$, forms a basis of $u.$ So $dim(u)=\rank(T-\lambda I|_{K_\lambda})=|\gamma|.$
        \item For each $i$, take the initial vector $w_i$ from $\beta_i.$ By basis extension theorem, the set $\{w_1,w_2,\cdots,w_q,u_1,u_2,\cdots,u_s\}$ forms a basis in $\ker(T-\lambda I).$
        Since $dim(K_\lambda)=\nullity(T-\lambda I|_{K_\lambda})+\rank(T-\lambda I|_{K_\lambda})=\nullity(T-\lambda I)+dim(u)=q+s+|\gamma|.$
        \item On the other hand, $\beta_i\subseteq im(T-\lambda I|_{K_\lambda})\subseteq im(T-\lambda I),$ for each $i,$ there exists a pre-image $v_i$ of the end vector of $\beta_i$. Let $\tilde{\beta}_i = \beta_i\cup\{v_i\}$ and the set $\tilde{\gamma} = \tilde{\beta_1}\cup\tilde{\beta_2}\cup\cdots\cup\tilde{\beta_q}\cup\{u_1\}\cup\{u_2\}\cup\cdots\cup\{u_s\}.$
        The initial vectors are linearly independent, so $\tilde{\gamma}$ is linearly independent, and $|\tilde{\gamma}|=|\gamma|+q+s.$ Hence, $\tilde{\gamma}$ forms a maximal linearly independent set in $K_\lambda$ and hence a basis.
    \end{enumerate}
\end{pfbox}


\begin{theorem}
    Let $\lambda$ and $\mu$ be distinct eigenvalues of $T:V\to V.$
    \begin{enumerate}
        \item $K_\lambda$ is a $T-$invariant subspace.
        \item The restriction of $T-\mu I$ to $K_\lambda$ is invertible.
    \end{enumerate}
\label{thm 7.7}
\end{theorem}

\begin{pfbox}
    For 1, use $T$ and $I$ can commute. Given that $T-\mu I|_{K_\lambda}$ is well-defined by 1, we want to show this restriction is bijective.
    By TFAE, taking any vector $\mathbf{v}\in \Ker(T-\mu I)\cap K_\lambda$, we hope $\mathbf{v}=\mathbf{0}.$ Suppose it does not. Since $\mathbf{v}\in K_\lambda$, there exists a smallest integer $p$ s.t. $(T-\lambda I)^p\mathbf{v}=\mathbf{0}.$ Let $\mathbf{y}=(T-\lambda I)^{p-1}\mathbf{v}$ such that $T\mathbf{y}=\lambda\mathbf{y}$.
    On the other hand, $\mathbf{v}\in \Ker(T-\mu I),$ we can show that $\mathbf{y}\in\Ker(T-\mu I),$ using commutativity btw $T$ and $I$. Then $T\mathbf{y}=\mu\mathbf{y},$ which yields a contradiction. 
\end{pfbox}

\begin{theorem}
    Let $a$ be a algebraic multiplicity of the eigenvalue $\lambda$ of $T$, then we have $K_\lambda=\Ker(T-\lambda I)^a.$
    \label{thm 7.8}
\end{theorem}

\begin{pfbox}
    $(\supseteq)$ is clear, for $(\subseteq)$, take $\mathbf{v}\in K_\lambda,$ write $p_T(x)=(x-\lambda)^a(x-\mu_1)^{b_1}(x-\mu_2)^{b_2}\cdots(x-\mu_k)^{b_k},$ where $\{\lambda, \mu_1, \mu_2, \cdots, \mu_k\}$ are distinct.
    By Cayley-Hamilton, one has: $p_T(T)(\mathbf{v})=(T-\lambda I)^a(T-\mu_1 I)^{b_1}(T-\mu_2 I)^{b_2}\cdots(T-\mu_k I)^{b_k}(\mathbf{v}) = \mathbf{O}(\mathbf{v})=\mathbf{0}.$
    Consider the map $(T-\mu_1 I)^{b_1}(T-\mu_2 I)^{b_2}\cdots(T-\mu_k I)^{b_k}$ restricted to $K_\lambda$, which is invertible by theorem~\eqref{thm 7.7}, then $(T-\lambda I)^a(\mathbf{v})=\mathbf{0}.$ This completes the proof. 
\end{pfbox}

\begin{remark}
    Recall theorem~\eqref{thm 7.4}, the above theorem asserts that $m\le a$.
\end{remark}

\begin{theorem}
    Let $\lambda$ and $\mu$ be distinct eigenvalues of $T$, then $K_\lambda\cap K_\mu=\{\mathbf{0}\}.$
\end{theorem}
\begin{pfbox}
    Take $\mathbf{v}\in K_\lambda\cap K_\mu$ and $p_T(x)=(x-\lambda)^a(x-\mu)^bq(x).$ Since $gcd\{(x-\lambda)^a, (x-\mu)^b\}=1,$ by Euclidean algorithm, there exists $f,g$ s.t. $f(x)(x-\lambda)^a+g(x)(x-\mu)^b=1$.
    Plug $x=T$ and apply the equation to $\mathbf{v},$ we must have $\mathbf{v}=\mathbf{0}$ by theorem~\eqref{thm 7.8} 
\end{pfbox}


\begin{theorem}[Primary Decomposition Theorem]
    Let $V$ be a finite-dimensional vector space over $\mathbb{C}$ and $T:V\to V$ be a linear operator. Let $\lambda_1,\lambda_2\cdots,\lambda_k$ be distinct eigenvalues of $T$, then:
    $$V=K_{\lambda_1}\oplus K_{\lambda_2}\oplus\cdots K_{\lambda_k}.$$
    \label{thm 7.10}
\end{theorem}

\begin{pfbox}
    Let $U=K_{\lambda_1}\oplus K_{\lambda_2}\oplus\cdots K_{\lambda_k}.$ We hope $dim(U)=dim(V).$
    (Direct) sum of space is a subspace, $U\subseteq V$ implies $dim(U)<dim(V).$
    Next, $dim(U)=\sum_{i=1}^{k}dim(K_{\lambda_i})\ge \sum_{i=1}^{k}\nullity(T-\lambda_i I)^{a_i}\ge\nullity(\prod_{i=1}^{k}(T-\lambda_i I)^{a_i})=\nullity(\mathbf{0})=k.$
    Combine $U\subseteq V$ and $dim(U)=dim(V),$ we conclude $U=V.$
\end{pfbox}

\begin{theorem}[JCF Theorem I]
    Let $V$ be a finite dimensional vector space over complex field, $T:V\to V$ be any linear operator, and the corresponding matrix representing $T$ w.r.t standard basis, $[T]^B_B$, 
    then there exists a basis $\gamma$ of $V$ formed by a union of Jordan chains, which is called a Jordan basis. Moreover, $[T]^B_B$ similar to its JCF, $[T]^\gamma_\gamma$.  
\end{theorem}

\begin{pfbox}
    By primary decomposition theorem~\eqref{thm 7.10} and theorem~\eqref{thm 7.6}, there exists such union of Jordan chains and forms a maximal linearly independent set and hence a basis.
\end{pfbox}

\begin{theorem}[JCF Theorem II \& III]
Let $A,B\in M_n(\mathbf{C})$, and up to a permutation of Jordan blocks, then:
    \begin{enumerate}
    \item $J_A$ is uniquely determined by $A.$
    \item $A\sim B$ if and only if $J_A=J_B.$
\end{enumerate}
\end{theorem}

\begin{pfbox}
    1 is trivial by definition. For 2, we focus on $(\Rightarrow)$. By similarity, $Spec(A)=Spec(B)$ and $(A-\lambda I)^i \sim (B-\lambda I)^i,$ which implies $\nullity(A-\lambda I)^i = \nullity(B-\lambda I)^i$ for each integer $i$. Hence, they have the same JCF. 
\end{pfbox}